%% ============================================================
%% preamble.tex
%% Shared preamble for
%%   Foundations of Sensitivity Analysis: Theory, Computation, and Applications
%%   Arriola & Hyman
%% ============================================================

%% ----------------------------------------------------------
%% PACKAGES
%% ----------------------------------------------------------
\usepackage[margin=0.75in]{geometry}
\usepackage{amsmath, amssymb, amsthm}
\usepackage{graphicx}
\graphicspath{{figures/}}
\usepackage{booktabs}
\usepackage{array}
\usepackage{longtable}
\usepackage{enumitem}
\usepackage{imakeidx}
\makeindex[columns=2, title=Index, intoc]
\usepackage{hyperref}
\usepackage[protrusion=true,expansion=false]{microtype}  %% improved microtypography
\hypersetup{
  colorlinks = true,
  linkcolor  = black,
  citecolor  = black,
  urlcolor   = black,
  pdftitle   = {Foundations of Sensitivity Analysis: Theory, Computation, and Applications},
  pdfauthor  = {Leon M. Arriola and James M. Hyman},
  pdfsubject = {Sensitivity Analysis, Uncertainty Quantification, Mathematical Modeling},
  pdfkeywords= {sensitivity analysis, adjoint methods, Sobol indices,
                SIR model, forward sensitivity, uncertainty quantification},
}
\usepackage{xcolor}
\usepackage{tcolorbox}
\usepackage{pgfplots}
\usepackage{tikz}
\usetikzlibrary{arrows.meta, positioning, decorations.markings, decorations.pathreplacing, calc, bending}
\pgfplotsset{compat=1.18}

%% ----------------------------------------------------------
%% THEOREM-STYLE ENVIRONMENTS
%% ----------------------------------------------------------
\theoremstyle{plain}
\newtheorem{theorem}{Theorem}[chapter]
\newtheorem{proposition}[theorem]{Proposition}
\newtheorem{lemma}[theorem]{Lemma}
\newtheorem{corollary}[theorem]{Corollary}

\theoremstyle{definition}
\newtheorem{definition}{Definition}[chapter]
\newtheorem{example}{Example}[chapter]
\newtheorem{remark}{Remark}[chapter]

%% ----------------------------------------------------------
%% BOOK-WIDE NOTATION
%% ----------------------------------------------------------

%% Sensitivity index S_p^q
\newcommand{\SI}[2]{S_{#1}^{#2}}

%% Parameter of Interest / Quantity of Interest
%% \xspace inserts a space after the acronym when followed by ordinary text,
%% but suppresses the space before punctuation and closing braces.
\usepackage{xspace}
\newcommand{\POI}{\textsc{poi}\xspace}
\newcommand{\QOI}{\textsc{qoi}\xspace}
\newcommand{\POIs}{\textsc{poi}s\xspace}
\newcommand{\QOIs}{\textsc{qoi}s\xspace}

%% Number sets
\newcommand{\R}{\mathbb{R}}
\newcommand{\N}{\mathbb{N}}

%% Probability / expectation
\newcommand{\E}{\mathbb{E}}
\newcommand{\Var}{\operatorname{Var}}

%% Derivatives
\newcommand{\pd}[2]{\frac{\partial #1}{\partial #2}}
\newcommand{\dd}[2]{\frac{d #1}{d #2}}

%% FSE shorthand
\newcommand{\FSE}{\textsc{fse}}

%% Unifying idea callback (inline, unnumbered)
\newcommand{\unifyingidea}[1]{\noindent\textit{Unifying Idea~#1.}\enspace}

%% Absolute / relative / regularized change
\newcommand{\Dabs}{\Delta u_{\mathrm{Abs}}}
\newcommand{\Drel}{\Delta u_{\mathrm{Rel}}}
\newcommand{\Dreg}{\Delta u_{\mathrm{Reg}}}

%% Math operators
\DeclareMathOperator{\arcsinh}{arcsinh}
\DeclareMathOperator{\sgn}{sgn}

%% ----------------------------------------------------------
%% PEDAGOGICAL ENVIRONMENTS
%% ----------------------------------------------------------

%% Self-check question with indented answer
\newenvironment{selfcheck}{%
  \medskip\noindent\textit{Self-check.}\enspace
}{\par}

\newenvironment{scAnswer}{%
  \nopagebreak\smallskip
  \noindent\hspace*{1.8em}%
  \begin{minipage}{0.90\textwidth}
    \footnotesize\textit{Answer:}\enspace
  }{\end{minipage}\medskip\par}

%% Numbered exercise environment (resets each chapter)
\newcounter{exercise}[chapter]
\renewcommand{\theexercise}{\thechapter.\arabic{exercise}}

\newenvironment{exercise}[1][]{%
  \refstepcounter{exercise}%
  \medskip\noindent\textbf{Exercise~\theexercise}%
  \ifx\relax#1\relax\else\ \textit{(#1)}\fi.\enspace
}{\par\medskip}

%% Bloom level tag (suppress in production by redefining to empty)
\newcommand{\bloom}[1]{\textnormal{\tiny\rmfamily[Bloom~#1]}}

%% ----------------------------------------------------------
%% CHAPTER-OPENING STYLE (book class default used;
%% customize here if needed)
%% ----------------------------------------------------------

%% Part / chapter appearance tweaks can go here.
%% For now, book class defaults are used throughout.
